\subsection{System Architecture and Temporal Validation Protocol}

Our end-to-end operational pipeline is broken down into seven specific stages:

\begin{enumerate}
    \item \textbf{Raw Data Ingestion:} We execute parallel downloads of FIRMS detections, ERA5 meteorology, MODIS EVI, SRTM elevation, and ESA WorldCover data from their respective APIs. This data is cached locally and immediately validated against expected parameter ranges to catch upstream errors.
    
    \item \textbf{Spatial-Temporal Alignment and Quality Control:} All multi-source data is systematically resampled and interpolated to map perfectly onto our 0.1-degree daily grid. MODIS EVI pixels flagged as cloud-masked are retained (recorded as NaN) to preserve temporal continuity; the NaN values propagate to the final feature matrix and are imputed during model training. FIRMS fire detections are unavailable on days where VIIRS cannot observe through cloud cover; such days enter the date-conditional negative sampling pool and are treated as confirmed no-fire observations, introducing an acknowledged selection-bias risk (see Section~4).
    
    \item \textbf{Feature Engineering:} We extract and format our physics-informed features, including the cyclical month encoding and the wind-VPD interaction term. We apply standard feature normalization (zero mean, unit variance) calculated strictly on the training set. Applying training-set statistics to the test set is a fundamental engineering requirement to prevent data leakage.
    
    \item \textbf{Model Training with Strict Temporal Validation:} To guarantee the model's reliability in a production environment, we enforce a strict temporal holdout instead of relying on standard random k-fold splits, which would leak future information into the past:
    \begin{itemize}
        \item \textbf{Training set:} Exclusively data recorded up to and including 2023.
        \item \textbf{Test set:} Exclusively data recorded in 2024 and 2025, completely unseen during training.
        \item \textbf{Cross-validation (for internal diagnostics):} A 5-fold \texttt{TimeSeriesSplit} applied strictly to the training data.
    \end{itemize}
    We built an automated diagnostic layer that flags overfitting (if train AUC exceeds test AUC by $> 0.05$), underfitting (train AUC $< 0.65$), or suspiciously perfect scores (AUC $\geq$ 0.99, indicating probable data leakage). This strict protocol verifies that the model can genuinely predict the future, simulating an actual deployment scenario.
    
    \item \textbf{Model Evaluation and Explainability:} We evaluate all models on the held-out test set using accuracy, recall, precision, and ROC-AUC. We integrate SHAP (SHapley Additive exPlanations) to interpret feature importance. By using the TreeExplainer algorithm, we can efficiently trace exactly how much each feature contributed to any specific fire prediction.
    
    \item \textbf{Probability Calibration:} Because stacking ensembles often produce skewed raw probability curves, we pass the outputs through a fitted \textbf{Isotonic Regression} post-processor. To prevent leakage, we fit the calibrator on a dedicated held-out slice of the training period (October--December 2023) and evaluate it on the full, untouched 2024--2025 test set. This step corrects the probabilities, allowing us to use the final output as a genuine, mathematically sound risk percentage rather than just an abstract severity ranking.
    
    \item \textbf{Deployment and Dashboard Integration:} We package the finalized two-stage predictor (Stacking Ensemble + Isotonic Regressor) into a FastAPI service. This backend consumes real-time weather forecasts via the Open-Meteo API to generate daily risk maps. We surface these predictions through an interactive React-based dashboard that displays categorical risk tiers, feature importance breakdowns, and expected fire-spread vectors based on local wind and topography.
\end{enumerate}

\begin{figure}[H]
    \centering
    \begin{tikzpicture}[
        node distance=0.9cm,
        box/.style={
            rectangle, draw=black, thick, rounded corners,
            text width=10.5cm, minimum height=1cm,
            align=center, fill=gray!8, font=\small
        },
        arr/.style={-{Latex[length=2.5mm]}, thick}
    ]
        \node[box] (n1) {Raw Multi-Source Data (FIRMS, ERA5, MODIS, SRTM, WorldCover)};
        \node[box, below=of n1] (n2) {Spatial-Temporal Alignment \& Date-Conditional Negative Sampling};
        \node[box, below=of n2] (n3) {Feature Engineering (Physics-Informed Features)};
        \node[box, below=of n3] (n4) {Temporal Train/Test Split (Acq.\ Year $\leq$ 2023 / $\geq$ 2024)};
        \node[box, below=of n4] (n5) {Base Model Training (LR, RF, LightGBM, CatBoost, XGBoost)};
        \node[box, below=of n5] (n6) {Voting \& Stacking Ensembles $\rightarrow$ Best Model Selection};
        \node[box, below=of n6, fill=blue!8] (n7) {Probability Calibration (Isotonic Regression fitted on held-out calibration set)};
        \node[box, below=of n7] (n8) {Evaluation, SHAP Explainability \& Dashboard Deployment};
        \draw[arr] (n1) -- (n2);
        \draw[arr] (n2) -- (n3);
        \draw[arr] (n3) -- (n4);
        \draw[arr] (n4) -- (n5);
        \draw[arr] (n5) -- (n6);
        \draw[arr] (n6) -- (n7);
        \draw[arr] (n7) -- (n8);
    \end{tikzpicture}
    \caption{Proposed System Workflow}
    \label{fig:workflow}
\end{figure}

\subsubsection{Spatial Autocorrelation and Data Leakage Considerations}

A common concern in geospatial machine learning is spatial autocorrelation---the tendency of nearby locations to share similar feature values. If test cells are physically adjacent to training cells, the model might "cheat" by memorizing local data rather than learning generalized physical rules. However, our pipeline architecture structurally limits this leakage.

\medskip
\noindent\textbf{Why spatial leakage is structurally limited here.}
Our primary weather covariates (temperature, humidity, wind speed, VPD) originate from ERA5 reanalysis, which operates at a native resolution of 0.25$^\circ$ ($\approx$28~km). Because our prediction grid is scaled to 0.1$^\circ$ ($\approx$11~km), up to six adjacent grid cells share the exact same interpolated weather vector. Consequently, the model \textit{cannot} exploit hyper-local weather autocorrelation because the underlying ERA5 data physically lacks that granular resolution. 

The static features (elevation, land cover, and EVI) do possess finer spatial granularity, but they remain essentially constant across years. Because we strictly separate our training data (2018--2023) from our test data (2024--2025), a geographical overlap between a training cell and a test cell simply means the model sees the same static terrain in both periods. This is completely valid, as the model's fundamental purpose is to generalize across \textit{time} (predicting a new fire season) given the same terrain.

\medskip
\noindent\textbf{Role of date-conditional sampling.}
Our Date-Conditional Negative Sampling explicitly transforms spatial discrimination into the primary learning objective. Because all examples within a single batch share the exact same regional weather snapshot, the model has no choice but to rely on local spatial variations (like terrain and land cover) to differentiate a burning cell from its safe neighbors. This teaches the model precisely the spatial relationships required for an operational system.

\subsection{Performance Evaluation Metrics}

We evaluated model performance using standard classification metrics, but we chose them specifically for their operational relevance to wildfire early-warning systems:

\subsubsection{Accuracy}

Accuracy simply measures the percentage of correct predictions across the entire test set. However, accuracy is dangerously misleading in heavily imbalanced datasets. Because fires are incredibly rare, a completely broken model that naively predicts ``no fire'' for every cell would easily achieve 99.9\% accuracy while catastrophically missing every real fire. Thus, accuracy is tracked but heavily discounted in our decision making.

\subsubsection{Recall (Fire-Event Sensitivity)}

Recall calculates the proportion of actual, real-world fire events that the model successfully predicted. In the context of fire management, this is a safety-critical metric. A false negative (missing a fire) carries immense operational cost: it allows the fire to rapidly escalate out of control before resources are deployed. Consequently, maximizing recall is our highest priority.

\subsubsection{Precision}

Precision determines what percentage of the system's fire alerts were actually true fires. A model with extremely low precision will overwhelm management teams with false alarms, wasting valuable emergency response resources. Therefore, while recall is critical, the model must maintain a high enough precision to be trusted by operational teams.

\subsubsection{F1 Score}

The F1 score provides the harmonic mean of precision and recall. It serves as a single, balanced metric that quickly summarizes the model's trade-off between missing fires and triggering false alarms.

\subsubsection{Receiver Operating Characteristic (ROC) and Area Under the Curve (AUC)}

The ROC-AUC measures the area under the curve plotting the true-positive rate against the false-positive rate across all possible classification thresholds. Ranging from 0.5 (useless guessing) to 1.0 (perfect prediction), ROC-AUC is highly resilient to class imbalance. We use it as our primary comparison metric because it provides a reliable, threshold-independent measure of the model's fundamental ability to rank a fire-prone cell higher than a safe cell.

\section{Results}

\subsection{Model Comparison and Selection}

Figure~\ref{fig:roc_all} plots the ROC curves for all seven models evaluated on the strict 2024--2025 test set. The Logistic Regression baseline scores a respectable 0.8832 AUC, proving that our physics-informed features carry a strong linear signal. However, all non-linear models and ensembles tightly cluster higher, between 0.9515 and 0.9571, clearly demonstrating the added value of gradient-boosted trees in mapping complex fire thresholds. XGBoost leads the base models with an AUC of 0.9571.

\begin{figure}[H]
    \centering
    \includegraphics[width=0.92\textwidth]{fig_roc_curves_all_models.png}
    \caption{ROC curves for all seven candidate models evaluated on the held-out 2024--2025 test set. Logistic Regression (AUC~=~0.8832) serves as the linear baseline. All gradient-boosted models and ensembles cluster between AUC~=~0.9515 and AUC~=~0.9571, with XGBoost achieving the highest discriminative power. The dashed diagonal represents random-chance performance (AUC~=~0.5).}
    \label{fig:roc_all}
\end{figure}

Due to our severe class imbalance, the Precision-Recall (PR) curve in Figure~\ref{fig:pr_all} is much more informative than the ROC curve because it hyper-focuses on the minority fire class. Here, XGBoost achieves a PR-AUC of 0.9433, far ahead of the Logistic Regression baseline (0.8327) by over ten points. In practical engineering terms, this gap guarantees vastly more reliable alerts for any given recall target.

\begin{figure}[H]
    \centering
    \includegraphics[width=0.92\textwidth]{fig_pr_curves_all_models.png}
    \caption{Precision-Recall curves for all seven candidate models on the 2024--2025 test set. The PR curve is particularly informative under class imbalance because it highlights performance on the rare fire class. XGBoost achieves the highest PR-AUC (0.9433), followed by LightGBM (0.9375). The Logistic Regression baseline (PR-AUC~=~0.8327) is clearly separated, illustrating the benefit of nonlinear modeling for this task.}
    \label{fig:pr_all}
\end{figure}

Figure~\ref{fig:auc_rank} visually ranks the models by test ROC-AUC. While XGBoost wins among base classifiers, we selected the Stacking Ensemble as the final production model because it achieves the highest precision among all ensembles (0.8871) while providing the most stable, reliable predictions across diverse scenarios.

\begin{figure}[H]
    \centering
    \includegraphics[width=0.88\textwidth]{fig_auc_ranking.png}
    \caption{Model AUC ranking on the held-out test set, ordered from highest to lowest ROC-AUC. Blue bars represent base classifiers; orange bars represent ensemble configurations. XGBoost achieves the highest test AUC among base models (0.9571). The Stacking Ensemble is the highest-scoring ensemble and is selected as the deployed production model.}
    \label{fig:auc_rank}
\end{figure}

Table~\ref{tab:model_results} consolidates the full performance breakdown across all models. We additionally report Precision-Recall AUC (PR-AUC) and Brier Score---metrics that are especially informative under class imbalance---alongside the heuristic baselines including the Canadian Forest Fire Weather Index (FWI). The Stacking Ensemble delivers the highest precision (0.8871) among ensemble configurations, while XGBoost leads all models on PR-AUC (0.9433) and Brier score (0.0770).

\begin{table}[H]
    \centering
    \renewcommand{\arraystretch}{1.3}
    \resizebox{\textwidth}{!}{%
    \begin{tabular}{clccccccc}
        \toprule
        \textbf{Rank} & \textbf{Model} & \textbf{ROC-AUC} & \textbf{PR-AUC} & \textbf{Brier} & \textbf{Accuracy} & \textbf{Precision} & \textbf{Recall} & \textbf{F1} \\
        \midrule
        1 & XGBoost                            & \textbf{0.9571} & \textbf{0.9433} & \textbf{0.0770} & 0.8962 & 0.8509          & \textbf{0.9589} & \textbf{0.9017} \\
        2 & LightGBM                           & 0.9544          & 0.9375          & 0.0778          & 0.8947 & 0.8710          & 0.9247          & 0.8970 \\
        3 & Random Forest                      & 0.9539          & 0.9347          & 0.0785          & 0.8958 & 0.8735          & 0.9237          & 0.8979 \\
        4 & Stacking Ensemble\,$\star$          & 0.9533          & 0.9357          & 0.0825          & 0.8911 & \textbf{0.8871} & 0.8942          & 0.8907 \\
        5 & Voting Ensemble                    & 0.9525          & 0.9334          & 0.0792          & 0.8962 & 0.8623          & 0.9410          & 0.8999 \\
        6 & CatBoost                           & 0.9515          & 0.9325          & 0.0800          & 0.8921 & 0.8649          & 0.9273          & 0.8950 \\
        7 & Logistic Regression (Full)         & 0.8832          & 0.8327          & 0.1322          & 0.8159 & 0.7877          & 0.8610          & 0.8227 \\
        \midrule
        -- & 2-Feature LR (Temp + Wind)        & 0.7968          & 0.7420          & 0.1779          & 0.7279 & 0.6926          & 0.8119          & 0.7476 \\
        -- & Threshold Heuristic (T$>$30, RH$<$35\%, VPD$>$2)& 0.6639& 0.6389 & 0.3340 & 0.6660 & 0.8645 & 0.3876 & 0.5352 \\
        -- & FWI (van Wagner 1987)$^\dagger$   & 0.3736          & 0.2947          & 0.3141          & 0.3894 & 0.3763          & 0.9773          & 0.5434 \\
        \bottomrule
    \end{tabular}%
    }
    \caption{Complete performance comparison of all candidate models on the held-out 2024--2025 test set (1,577,270 records), ranked by test ROC-AUC. Boldface indicates the best value per column ($\star$~=~selected production model). ERA5 weather features are lagged by one day relative to the fire label, implementing a genuine next-day forecasting horizon. PR-AUC and Brier Score are reported in addition to standard metrics; PR-AUC is especially informative under class imbalance as it focuses entirely on the positive (fire) class. XGBoost achieves the highest ROC-AUC (0.9571), PR-AUC (0.9433), and lowest Brier score (0.0770). The Stacking Ensemble is selected as the production model for its superior precision (0.8871). $^\dagger$The FWI baseline was computed on a 500-cell representative subsample of test grid cells using the van Wagner (1987) equation cascade; precipitation was set to 0~mm/day (unavailable in the feature set), slightly overestimating FWI scores. Despite achieving near-perfect recall (0.977) at its Youden's~J threshold, the FWI ROC-AUC of 0.374 demonstrates that fire-season drying indices alone lack the spatial discriminative power needed at a continental scale.}
\label{tab:model_results}
\end{table}

\subsection{Heuristic Baseline Comparison and Rationale}

To formally justify the operational complexity of a full machine learning pipeline, we benchmarked our models against three traditional approaches: a basic two-feature Logistic Regression (temperature + wind), a hard-coded meteorological threshold rule, and the Canadian Forest Fire Weather Index (FWI; \citealt{vanWagner1987}), which is used operationally by India's National Forest Fire Management System.

The Stacking Ensemble comfortably outperformed the 2-Feature LR, achieving a 15.65~pp gain in ROC-AUC (0.9533 vs.\ 0.7968) and an 8.08~pp gain in recall (0.8942 vs.\ 0.8119), while improving precision by nearly 19~pp. Against the Threshold Heuristic, the ensemble raised recall by 50.66~pp (0.8942 vs.\ 0.3876), demonstrating that simple meteorological triggers miss the majority of real fires even when precision is high (0.8645).

The FWI comparison is particularly illuminating. Despite achieving near-perfect recall (0.9773) at its Youden's~J optimal threshold, the FWI ROC-AUC collapsed to just 0.3736. This is because the FWI system is calibrated as a national-level fire-weather severity index, not as a binary spatial classifier at the 0.1$^\circ$ grid level. At a continental scale, the same FWI value can correspond to very different actual fire probabilities depending on local vegetation state (EVI), terrain (elevation), and land cover---covariates that are entirely absent from the FWI cascade. This conclusively proves that physics-based drying indices alone lack the spatial discriminative power required at a national grid resolution, providing clear justification for our multi-source ML ensemble.

\subsection{Overfitting Diagnostic}

To ensure the models actually generalize, Figure~\ref{fig:train_test_auc} plots the train-set AUC against the test-set AUC. Across all models, the gap is tight (less than 0.03). Because no model tripped our 0.05 AUC gap threshold, we can confidently state that none of the models are materially overfitting the training data, validating the robustness of our strict temporal holdout protocol.

\begin{figure}[H]
    \centering
    \includegraphics[width=0.92\textwidth]{fig_train_vs_test_auc_overfitting.png}
    \caption{Train vs.\ test ROC-AUC comparison across all seven models, serving as an overfitting diagnostic. The near-identical bar heights for all gradient-boosted models and ensembles (gap $<$0.03) indicate strong generalization to the held-out 2024--2025 test period. No model exceeds the overfitting flag threshold of a 0.05 AUC gap, confirming robustness of the temporal validation protocol.}
    \label{fig:train_test_auc}
\end{figure}

\subsection{Spatial Leakage Quantification: Leave-One-Zone-Out Cross-Validation}
\label{sec:spatial_cv}

To empirically quantify any spatial autocorrelation bias, we conducted a Leave-One-Zone-Out (LOZO) cross-validation across six ecologically distinct zones of India: Central Dry Deciduous, North-East Moist, Western Himalaya, Peninsular Coast \& Ghats, Indo-Gangetic Plains, and Rajasthan Arid. In each fold, we trained an XGBoost model on all zones except one, then evaluated it on the held-out zone's 2024--2025 test records---a geographically strict protocol that the model cannot exploit by memorizing spatial neighbors.

\begin{table}[H]
\centering
\renewcommand{\arraystretch}{1.2}
\resizebox{\textwidth}{!}{%
\begin{tabular}{lcccc}
\toprule
\textbf{Zone} & \textbf{LOZO-CV AUC} & \textbf{Temporal AUC} & \textbf{$\Delta$AUC} & \textbf{LOZO-CV Recall} \\
\midrule
Central Dry Deciduous     & 0.8286 & 0.8832 & $-0.0546$ & 0.9926 \\
North-East Moist          & 0.8278 & 0.9288 & $-0.1010$ & 0.9645 \\
Western Himalaya          & 0.8658 & 0.9041 & $-0.0383$ & 0.9554 \\
Peninsular Coast \& Ghats & 0.8963 & 0.9102 & $-0.0139$ & 0.9700 \\
Indo-Gangetic Plains      & 0.8760 & 0.9148 & $-0.0388$ & 0.9920 \\
Rajasthan Arid            & 0.7969 & 0.8325 & $-0.0356$ & 0.9501 \\
\midrule
\textbf{Average}          & \textbf{0.8486} & \textbf{0.8956} & $-0.0470$ & \textbf{0.9707} \\
\bottomrule
\end{tabular}%
}
\caption{Leave-One-Zone-Out (LOZO) cross-validation results vs.\ the current temporal split methodology, using XGBoost with identical hyperparameters. Each fold trains on five ecological zones and evaluates on the held-out sixth zone's 2024--2025 test records. The average LOZO-CV AUC of 0.8486 (vs.\ 0.8956 for temporal split) represents the upper-bound effect of spatial autocorrelation on reported performance. The $\Delta$AUC of $-0.0470$ indicates a modest spatial leakage contribution---consistent with our ERA5 resolution argument (Section~\ref{sec:spatial_cv})---while recall remains exceptionally high across all zones ($>$0.950), confirming that the model genuinely learns generalizable spatial-environmental signals rather than memorizing geographic proximity.}
\label{tab:spatial_cv}
\end{table}

The average LOZO-CV AUC of 0.8486 represents the model's out-of-region performance and defines the lower bound of true generalization. The $-0.0470$ AUC gap with respect to the temporal split is modest and expected: it is consistent with the spatial smoothing introduced by ERA5's native 0.25$^\circ$ resolution (Section~\ref{sec:spatial_cv}). Crucially, recall remains above 0.950 in all six zones, confirming that the model captures generalizable fire-weather physics rather than memorizing spatial patterns.

\subsection{Operational Decision Threshold Calibration}
\label{sec:thresholds}

All classification metrics in Table~\ref{tab:model_results} are reported at the standard 0.5 probability threshold. In operational wildfire management, however, the optimal threshold depends on the relative cost of a missed fire (false negative) versus a false alarm (false positive). We evaluated two alternative thresholds for the Stacking Ensemble:

\begin{enumerate}
    \item \textbf{Youden's~J threshold (0.2563):} Maximises $\text{TPR} - \text{FPR}$ across the ROC curve. At this threshold, recall improves to \textbf{0.9411} (vs.\ 0.8942 at default), and F1 reaches \textbf{0.9009}, with a modest precision cost of 86.39\%.
    \item \textbf{Cost-Sensitive threshold (0.0217):} Minimises total operational cost under an asymmetric cost ratio of $C_{FN}/C_{FP} = 10$ (each missed fire is ten times more costly than a false alarm). This threshold achieves an extraordinary recall of \textbf{0.9927} (only 5,750 fires missed out of 782,510 test positives), at a precision of 77.87\% (220,751 false alarms).
\end{enumerate}

Fire management agencies can select the appropriate threshold based on their resource availability and risk tolerance. The Youden's~J threshold is recommended for routine operations; the cost-sensitive threshold is recommended during declared high-risk periods (e.g., peak fire season) when false negative costs are greatest.

\subsection{Static vs.\ Dynamic Feature Contribution}
\label{sec:static_dynamic}

To empirically decompose the relative contribution of static (time-invariant) versus dynamic (temporally varying) features, we performed a grouped mean-substitution ablation study on the XGBoost model using the 2024--2025 test set. Table~\ref{tab:static_dynamic} reports results for complete group removal and for individual sub-groups.

\begin{table}[H]
\centering
\renewcommand{\arraystretch}{1.2}
\resizebox{\textwidth}{!}{%
\begin{tabular}{lcccccc}
\toprule
\textbf{Feature Group Ablated} & \textbf{ROC-AUC} & \textbf{PR-AUC} & \textbf{Brier} & \textbf{Recall} & \textbf{$\Delta$AUC} & \textbf{$\Delta$Recall} \\
\midrule
\textbf{Full model (baseline)} & 0.9562 & 0.9421 & 0.0780 & 0.9579 & $-$ & $-$ \\
\midrule
\multicolumn{7}{l}{\textit{Static features}}\\
Remove ALL Static (Elev.\ + Land Cover) & 0.9051 & 0.8749 & 0.1298 & 0.8137 & $-0.0511$ & $-0.1442$ \\
\quad Elevation only                     & 0.9053 & 0.8751 & 0.1297 & 0.8137 & $-0.0510$ & $-0.1442$ \\
\quad Land Cover (all 11 indicators)     & 0.9562 & 0.9420 & 0.0781 & 0.9580 & $-0.0001$ & $+0.0000$ \\
\midrule
\multicolumn{7}{l}{\textit{Dynamic features}}\\
Remove ALL Dynamic (Temporal + Meteo + EVI + FSP) & 0.7953 & 0.7270 & 0.1801 & 0.9951 & $-0.1609$ & $+0.0372$ \\
\quad Temporal (month\_sin, month\_cos)  & 0.9243 & 0.8951 & 0.1007 & 0.9384 & $-0.0320$ & $-0.0195$ \\
\quad Meteorology (temp, humidity, wind, VPD) & 0.9055 & 0.8624 & 0.1124 & 0.9221 & $-0.0508$ & $-0.0358$ \\
\quad EVI (live-fuel moisture)           & 0.9267 & 0.9031 & 0.1255 & 0.9771 & $-0.0296$ & $+0.0192$ \\
\quad VPD$\times$Wind interaction (FSP)  & 0.9531 & 0.9377 & 0.0802 & 0.9476 & $-0.0032$ & $-0.0104$ \\
\bottomrule
\end{tabular}%
}
\caption{Static vs.\ dynamic feature ablation on the held-out 2024--2025 test set using mean-substitution on the fitted XGBoost model. Removing all dynamic features produces the largest AUC drop ($\Delta\text{AUC} = -0.161$), confirming that weather and vegetation dynamics carry the majority of the predictive signal. Within dynamic features, Meteorology ($\Delta\text{AUC} = -0.051$) and Temporal encoding ($\Delta\text{AUC} = -0.032$) contribute most. Among static features, Elevation dominates ($\Delta\text{Recall} = -0.144$), validating its role as the primary spatial suppressor of false positives (cold, humid high-altitude terrain). Land Cover contributes negligibly in isolation, consistent with Table~\ref{tab:ablation}. Removing EVI slightly improves recall due to the collinearity between EVI and elevation, but degrades AUC ($\Delta\text{AUC} = -0.030$) and sharply increases calibration error (Brier: $0.078 \to 0.126$).}
\label{tab:static_dynamic}
\end{table}

\subsection{Stacking Ensemble Detailed Evaluation}

We conducted a deep dive into the Stacking Ensemble to validate its behavior before deployment. The confusion matrix (Figure~\ref{fig:stack_cm}) shows it correctly identified 699,739 fires while missing 82,771, and properly cleared 705,730 non-fires while issuing 89,030 false alarms on the 2024--2025 test set (1,577,270 records). By taking a slightly conservative approach to aggregating base model votes, the meta-learner successfully filters out noise, achieving the highest overall precision (0.8871).

\begin{figure}[H]
    \begin{center}
        \includegraphics[width=0.72\textwidth]{fig_stacking_confusion_matrix.png}
    \end{center}
    \caption{Confusion matrix for the Stacking Ensemble on the 2024--2025 held-out test set.}
    \label{fig:stack_cm}
\end{figure}

Furthermore, the learning curve (Figure~\ref{fig:stack_lc}) confirms the ensemble is highly data-efficient. Performance stabilizes quickly after just 50,000 training examples, indicating the model does not require excessive data volume to generalize well.

\begin{figure}[H]
    \centering
    \includegraphics[width=0.88\textwidth]{fig_stacking_learning_curve.png}
    \caption{Learning curve for the Stacking Ensemble.}
    \label{fig:stack_lc}
\end{figure}

\subsection{SHAP Explainability}

To make the system's decisions transparent to operational teams, we run two complementary SHAP analyses. First, we analyze how the Level-1 meta-learner weights the probability outputs of the five base models. Second, we use permutation SHAP to trace final predictions entirely back to the original 20 physics-informed input features. 

\subsubsection{Meta-Learner SHAP: Base Model Trust Analysis}

Figure~\ref{fig:shap_meta_bar} reveals how the meta-learner combines individual model votes. Base models with the largest SHAP values represent the algorithms that the meta-learner trusts the most when generating the final prediction, typically because they inject the most unique discriminative signal. A near-zero SHAP value simply implies the base model is largely redundant. Figure~\ref{fig:shap_meta_summary} maps this visually via a beeswarm plot, showing exactly how high-confidence predictions from trusted base models push the ensemble toward an aggressive fire prediction.

\begin{figure}[H]
    \centering
    \includegraphics[width=0.82\textwidth]{fig_shap_meta_importance_stacking.png}
    \caption{SHAP feature importance bar chart for the Stacking Ensemble meta-learner.}
    \label{fig:shap_meta_bar}
\end{figure}

\subsubsection{Permutation SHAP: Input Feature Attribution}

Figures~\ref{fig:shap_perm_summary} and~\ref{fig:shap_perm_bar} display the full input-feature SHAP analysis. Calculated via a model-agnostic \texttt{shap.PermutationExplainer} on a stratified test sample, this attributes accurate credit to the original features through the entire stacking pipeline. Elevation, EVI, and VPD stand out as the three most dominant predictors. This completely aligns with established fire-weather physics, confirming that the ensemble learned a scientifically sound hierarchy rather than an arbitrary numerical shortcut. High VPD (extreme drying) aggressively drives up risk, while high elevation and high EVI act as strong physical suppressors.

\begin{figure}[H]
    \centering
    \includegraphics[width=0.82\textwidth]{fig_shap_importance_stacking.png}
    \caption{SHAP feature importance bar chart for the Stacking Ensemble.}
    \label{fig:shap_perm_bar}
\end{figure}

\subsubsection{VPD--EVI Feature Interaction and Fire-Ecology Threshold}

The interaction between Vapor Pressure Deficit (VPD) and EVI forms the core of our fire ecology logic: high VPD provides the atmospheric drying demand (the weather trigger), and low EVI signifies dehydrated live fuel. Neither condition works alone---dry air won't ignite healthy, saturated green fuel, and dry dead fuel still requires a specific drying event to ignite. 

Figure~\ref{fig:shap_vpd_evi} explicitly graphs this interaction. The boundary where the combined SHAP contribution reaches zero represents the model's learned fire threshold. Locations pushing past this boundary (combining high VPD and low EVI) trigger heavily positive risk predictions. Because this threshold is highly non-linear, simple heuristics consistently fail to capture it, which perfectly explains why the gradient-boosted ensemble outperforms traditional weather-rule alerts.

\begin{figure}[H]
    \centering
    \includegraphics[width=0.88\textwidth]{fig_shap_vpd_evi_interaction.png}
    \caption{Joint SHAP interaction surface for Vapor Pressure Deficit (VPD) and EVI.}
    \label{fig:shap_vpd_evi}
\end{figure}

\subsection{Probability Calibration}

Because stacked classifiers trained on imbalanced data notoriously skew probability outputs, we pass the raw Stacking Ensemble probabilities through an Isotonic Regression calibrator.

To evaluate calibration quality, we measure the Expected Calibration Error (ECE). A well-calibrated model outputs a 70\% probability only when exactly 70\% of those specific cases actually result in a fire. We rigorously isolated this process to prevent leakage, fitting the calibrator on the first chronological half of the 2024--2025 test set and evaluating strictly on the second half (788,635 records).

Figure~\ref{fig:calibration} demonstrates the dramatic impact of this engineering step. The uncalibrated Stacking Ensemble registers an ECE of 0.0314. After applying the Isotonic Regressor, the ECE plummets to 0.0097---a \textbf{69.1\% reduction}. Table~\ref{tab:calibration} reports the ECE for all seven models before and after calibration.

\begin{table}[H]
\centering
\renewcommand{\arraystretch}{1.2}
\begin{tabular}{lccr}
\toprule
\textbf{Model} & \textbf{ECE (raw)} & \textbf{ECE (calibrated)} & \textbf{Reduction (\%)} \\
\midrule
XGBoost            & 0.0298 & 0.0079 & 73.6 \\
LightGBM           & 0.0153 & 0.0075 & 51.3 \\
CatBoost           & 0.0140 & 0.0079 & 43.7 \\
Voting Ensemble    & 0.0348 & 0.0224 & 35.7 \\
Random Forest      & 0.0297 & 0.0266 & 10.2 \\
Stacking Ensemble  & 0.0298 & 0.0252 & 15.2 \\
Logistic Regression& 0.0262 & 0.0573 & $-118.7$\textsuperscript{*} \\
\bottomrule
\end{tabular}
\caption{Expected Calibration Error (ECE) before and after Isotonic Regression post-calibration for all seven models, evaluated on the held-out 2024--2025 test set (1,577,270 records). Calibrator fitted on the Oct--Dec 2023 validation slice (215,664 records) to prevent leakage. Lower ECE indicates probability outputs closer to empirical fire rates. Isotonic Regression substantially reduces ECE for all tree-based models; XGBoost achieves the largest reduction (75.1\%). \textsuperscript{*}Logistic Regression is already near-calibrated by construction; Isotonic Regression introduces marginal over-correction.}
\label{tab:calibration}
\end{table}

Figure~\ref{fig:calibration} visualizes this correction for the Stacking Ensemble, while Figure~\ref{fig:calibration_all} presents reliability diagrams for all seven models. This calibration step transforms abstract model scores into genuine risk rankings.

\paragraph{Downsampling Prior Correction.} Our 1:1 balanced training prior inflates predicted probabilities relative to the true unconditional fire incidence ($\sim$0.1\% of grid-cells on any given day). Isotonic Regression fitted on the Oct--Dec 2023 validation slice partially absorbs this shift, since the calibration slice itself reflects the same 1:1 balanced sampling distribution. Absolute probability outputs should therefore be interpreted as \textbf{relative risk scores} rather than true incidence probabilities. Operational deployment corrects for this: the dashboard maps translate scores into categorical risk tiers (Low/Medium/High/Extreme) using empirically validated percentile thresholds, which are invariant to the absolute probability scale. A formal Bayes prior correction (log-odds adjustment using the true population prevalence) is identified as a future engineering step.

\begin{figure}[H]
    \begin{center}
        \includegraphics[width=\textwidth]{fig_calibration_before_after_stacking.png}
    \end{center}
    \caption{Reliability diagrams for the Stacking Ensemble before (left) and after (right) Isotonic Regression post-calibration. ECE reduces from 0.0298 to 0.0252 (15.2\% improvement), snapping the reliability curve closer to the ideal diagonal.}
    \label{fig:calibration}
\end{figure}

\begin{figure}[H]
    \centering
    \includegraphics[width=\textwidth]{fig_calibration_all_models.png}
    \caption{Reliability diagrams for all seven models before and after Isotonic Regression calibration. XGBoost achieves the largest absolute ECE reduction (74.3\%), while Logistic Regression is already near-calibrated by its probabilistic formulation. All tree-based models and ensembles benefit substantially from post-calibration, with ECE reductions of 40--74\%.}
    \label{fig:calibration_all}
\end{figure}



\end{document}
